\section{符号说明}

% 符号定义维护于 docs/notations.md；增补或改名时同步更新两处。
% 使用 longtable 保留续表表头，后续新增条目无需压缩字号或行距。
本文主要符号及单位见表\ref{tab:notations}。
下标 $a$、$s$、$0$ 分别表示烘房空气、药材表面和初始状态，
正文中的 $T$、$T_a$ 和 $T_s$ 统一采用摄氏温度，单位为 ${}^\circ\mathrm{C}$。凡进入 Arrhenius 经验公式，单独定义 $T_K=T+273.15$，并以 $T_K$ 代入指数项；后续各问均遵循此约定。

\begingroup
% 与项目 tabular 表格沿用相同字号、行距，不修改公共模板。
\zihao{-5}
\renewcommand{\arraystretch}{1.38}
\begin{longtable}{@{}p{0.13\textwidth}p{0.60\textwidth}p{0.20\textwidth}@{}}
\caption{核心符号及单位}\label{tab:notations}\\
\toprule
符号 & 含义 & 单位 \\
\midrule
\endfirsthead
\multicolumn{3}{c}{\textbf{表\thetable\quad 核心符号及单位（续）}}\\
\toprule
符号 & 含义 & 单位 \\
\midrule
\endhead
\midrule
\multicolumn{3}{r}{续下页}\\
\endfoot
\bottomrule
\endlastfoot
$t$ & 烘干时间 & $\mathrm{s}$ \\
$t_e$ & 全域最大水分浓度首次降至阈值的临界时刻 & $\mathrm{s}$ \\
$\xi$ & 材料坐标，$\xi=r/R(t)$ & $1$ \\
$C_{\max}$ & 全域最大水分浓度 & $\mathrm{kg/kg}$ \\
$r$ & 到圆柱中心轴的径向距离 & $\mathrm{m}$ \\
$R$ & 药材半径；尺寸变化时记为 $R(t)$ & $\mathrm{m}$ \\
$L$ & 药材长度 & $\mathrm{m}$ \\
$T$ & 药材温度，正文采用摄氏温度 & ${}^\circ\mathrm{C}$ \\
$C$ & 药材水分浓度，即干基含水率 & $\mathrm{kg/kg}$ \\
$\rho$ & 药材密度 & $\mathrm{kg/m^3}$ \\
$c_p$ & 药材比热容 & $\mathrm{J/(kg\cdot K)}$ \\
$k$ & 热传导系数 & $\mathrm{W/(m\cdot K)}$ \\
$D$ & 水分扩散系数 & $\mathrm{m^2/s}$ \\
$h$ & 对流换热系数 & $\mathrm{W/(m^2\cdot K)}$ \\
$h_m$ & 对流传质系数 & $\mathrm{m/s}$ \\
$T_a$ & 烘房空气温度，正文采用摄氏温度 & ${}^\circ\mathrm{C}$ \\
$C_a$ & 烘房空气水分浓度，采用附件口径 & $\mathrm{kg/kg}$ \\
$T_s$ & 药材表面温度，正文采用摄氏温度 & ${}^\circ\mathrm{C}$ \\
$C_s$ & 药材表面水分浓度，即表面干基含水率 & $\mathrm{kg/kg}$ \\
$\alpha$ & 热扩散率，$\alpha=k/(\rho c_p)$ & $\mathrm{m^2/s}$ \\
$\tau_T$ & 径向热扩散特征时间 & $\mathrm{s}$ \\
$\tau_C$ & 径向水分扩散特征时间 & $\mathrm{s}$ \\
$N$ & 径向控制体总数 & $1$ \\
$\Delta r$ & 均匀径向网格宽度 & $\mathrm{m}$ \\
$i$ & 控制体或节点编号 & $1$ \\
\end{longtable}
\endgroup

表中单位“$1$”表示无量纲。
物性参数及数值离散符号均在相应问题中定义。
